\begin{figure}[htbp]
\centering
\begin{tikzpicture}[
    node distance  = 0.7cm,
    etapa/.style   = {
        rectangle, rounded corners = 5pt,
        draw = black!55, line width = 0.55pt,
        text width = 10cm, minimum height = 1.3cm,
        align = center, font = \small
    },
    flecha/.style  = {->, >= Stealth, line width = 1pt},
    nota/.style    = {font = \scriptsize\itshape, text = black!45}
]

\node[etapa, fill = gray!12] (e1) {
    \textbf{Etapa 1 — Recopilación de instancias}\\[3pt]
    \footnotesize Benchmarks estándar y librerias $\longrightarrow$ formato SDPA sparse
};

\node[etapa, fill = blue!7, below = of e1] (e2) {
    \textbf{Etapa 2 — Benchmarking de solvers}\\[3pt]
    \footnotesize Tiempo de cómputo $\cdot$ iteraciones $\cdot$ calidad de solución
};

\node[etapa, fill = blue!7, below = of e2] (e3) {
    \textbf{Etapa 3 — Extracción de features}\\[3pt]
    \footnotesize Estructura de bloques $\cdot$ densidad $\cdot$ conectividad $\cdot$ cordalidad
};

\node[etapa, fill = teal!10, below = of e3] (e4) {
    \textbf{Etapa 4 — Análisis ISA}\\[3pt]
    \footnotesize Generación de espacio de instancias $\longrightarrow$ regiones de ventaja comparativa de solvers
};

\node[etapa, fill = orange!12, below = of e4] (e5) {
    \textbf{Etapa 5 — Generación evolutiva}\\[3pt]
    \footnotesize Algoritmo Genético $\longrightarrow$ instancias sintéticas 
};

\draw[flecha] (e1) -- (e2);
\draw[flecha] (e2) -- (e3);
\draw[flecha] (e3) -- (e4);
\draw[flecha] (e4) -- (e5);

% Retroalimentación: E5 → E1
\draw[flecha, dashed, black!40]
    (e5.east) -- ++(1.8, 0) coordinate (aux) |- (e1.east);
\node[nota] at ($(aux)!0.5!(aux |- e3)$) [right = 3pt]
    {retroalimentación};

\end{tikzpicture}
\caption{Pipeline metodológico de la investigación.
         La flecha discontinua indica que las instancias generadas en la Etapa~5
         se incorporan al conjunto experimental y realimentan el análisis desde la Etapa~1.}
\label{fig:pipeline}
\end{figure}
